\chapter{Conclusiones y trabajo futuro}
\label{cap:conclusiones}

% =====================================================================================
% CAPITULO BLOQUEADO. Depende del capitulo 10.
% =====================================================================================

\section{Cumplimiento de los objetivos}
\label{sec:cumplimiento}
Repasar uno a uno los objetivos del capítulo~\ref{cap:intro} y decir, sin adornos, cuáles se
cumplieron y cuáles no.

\section{Aportaciones}
\label{sec:aportaciones}
% Las cuatro del capitulo 1, ya demostradas.

\section{Limitaciones}
\label{sec:limitaciones_finales}
% NO minimizarlas: declararlas hace creible el resto.
\begin{itemize}
    \item El modelo de ruido \textbf{no incluye crosstalk}: las conclusiones son válidas
          dentro de ese modelo.
    \item No hay validación en \textbf{hardware real}: el plan abierto de IBM da 10 minutos
          de QPU al mes.
    \item El \textit{ground truth} exige simulación, lo que \textbf{acota el tamaño} de los
          circuitos.
    \item La reproducibilidad es del plan, \textbf{no del bit}
          (capítulo~\ref{cap:rigor}).
\end{itemize}

\section{Trabajo futuro}
\label{sec:futuro}
% Fuente: TFM-Quantum/IDEAS_FUTURAS.md
Validación en hardware real, mitigación de errores de lectura, crosstalk en el modelo de
ruido, y ampliación de la batería de observables.
